\documentclass[]{ufpe-exercises}
\usepackage{hyperref}

\curso{Ciências Contábeis}
\discente{Pedro Malta Schuler Caloête}
\centro{Centro de Ciências Sociais Aplicadas}
\departamento{Departamento de Ciências Contábeis e Atuariais}
\titulo{Lista 2 - Exercício 3}
\disciplina{Contabilometria}
\docente{Giuseppe Trevisan}
\periodo{2026.1}
\corcabecalho{red4}

\begin{document}

\cabecalho

\begin{questions}

	\question Conforme Exercício 3, disponível no site do professor [\url{https://sites.google.com/view/giuseppetrevisan/teaching}]

	\begin{parts}

		\part

		\begin{solution}
			Como estamos trabalhando com um modelo \(log-lin\), precisamos linearizar nossa variável dependente antes de efetuarmos os cálculos. desejamos obter o seguinte modelo:

			\[log(Preco Acao) = \beta_{0} + \beta_{1} \times Gastos Dos Pares\]

			Podemos, então, montar nossa tabela:

			\begin{center}
				\begin{tblr}{
						colspec = {X[c] X[c] X[c] X[c] X[c] X[c]},
						hlines,
						vlines
					}
					\hline
					Empresa  & Preço da Ação (Y) & log(Preço da Ação) & Gastos dos Pares (X) & \(X \times Y\) & \(X^2\) \\
					1        & 21,0              & 3,045              & 35                   & 625            & 1.225   \\
					2        & 23,5              & 3,157              & 28                   & 784            & 784     \\
					3        & 14,6              & 2,681              & 60                   & 1225           & 3.600   \\
					4        & 13,1              & 2,573              & 71                   & 400            & 5.041   \\
					5        & 20,0              & 2,996              & 44                   & 625            & 1.936   \\
					6        & 18,1              & 2,896              & 45                   & 900            & 2.025   \\
					7        & 20,2              & 3,006              & 58                   & 1225           & 3.364   \\
					8        & 18,5              & 2,918              & 44                   & 625            & 1.936   \\
					\(\sum\) & ---               & 23,272             & 385                  & 1.103,398      & 19.911  \\
				\end{tblr}
			\end{center}

			\(\hat{\beta_{1}}\) pode ser definido da seguinte forma:
			\[\hat{\beta_{1}} = \frac{n\sum (X_{i}Y_{i}) - \sum Y_{i}\sum X_{i}}{n\sum X^2 - (\sum X_{i})^2}\]
			Agora podemos utilizar os valores da tabela criada para calcular \(\hat{\beta_{1}}\).
			\[ \hat{\beta_{1}} = \frac{8 \times 1103,398 - 385 \times 23,272}{8 \times 19911 - (385)^2} = -0,01198 \]
			\(\hat{\beta_{0}}\) pode ser definido como \(\hat{\beta_{0}} = \bar{Y} - \beta_{1} \times \bar{X}\), logo:
			\[ \hat{\beta_{0}} = \frac{23,272}{8} - 0,01198 \times \frac{385}{8} = 3,48554 \]

			Sabendo que nosso regressor é contínuo, e agora que temos os valores de \(\hat{\beta_{0}}\) e \(\hat{\beta_{1}}\) podemos interpretar o seguinte: quando os gastos com controle interno dos pares forem igual a 0, o preço da ação será, em média, igual a R\$ 3.485,54 mil reais. O \(\hat{\beta_{1}}\) nos informa que, quando os gatos com controle interno dos pares aumenta em mil reais, o preço da ação diminui, em média, em 1,198\%.

		\end{solution}

		\vspace{0.3cm}

		\part

		\begin{solution}

			A variância do erro, denotada por \(\sigma^{2}\), é um parâmetro populacional e, portanto, desconhecido. Para estimá-lo podemos utilizar os seguintes estimadores:
			\[
				S^{2} = \frac{\sum \hat{u_i^2}}{n-k} \quad
				\hat{\sigma^2}= \frac{\sum \hat{u_i^2}}{n}
			\]

			Inicialmente precisamos calcular o somatório do quadrado dos resíduos (SQR).
			\[SQR = \sum (Y - \widehat{Y})^{2}\]
			\[ [ 3,0 - (\hat{\beta_{0}} + \hat{\beta_{1}} \times 25)]^{2} =  0,08563 \]
			\[ [ 3,5 - (\hat{\beta_{0}} + \hat{\beta_{1}} \times 28)]^{2} =  0,00295 \]
			\[ [ 3,6 - (\hat{\beta_{0}} + \hat{\beta_{1}} \times 35)]^{2} =  0,04118 \]
			\[ [ 3,1 - (\hat{\beta_{0}} + \hat{\beta_{1}} \times 20)]^{2} =  0,00391 \]
			\[ [ 3,3 - (\hat{\beta_{0}} + \hat{\beta_{1}} \times 25)]^{2} =  0,00005 \]
			\[ [ 3,8 - (\hat{\beta_{0}} + \hat{\beta_{1}} \times 30)]^{2} =  0,06361 \]
			\[ [ 4,0 - (\hat{\beta_{0}} + \hat{\beta_{1}} \times 35)]^{2} =  0,03884 \]
			\[ [ 3,1 - (\hat{\beta_{0}} + \hat{\beta_{1}} \times 25)]^{2} =  0,03711 \]
			\[ [ 3,2 - (\hat{\beta_{0}} + \hat{\beta_{1}} \times 20)]^{2} =  0,02641 \]
			\[ [ 3,5 - (\hat{\beta_{0}} + \hat{\beta_{1}} \times 30)]^{2} =  0,00228 \]
			\[ SQR = 0,30197\]

			Portanto:
			\[
				S^{2} = \frac{0,30197}{10-2} = \frac{0,30197}{8} = 0,03775 \quad
				\hat{\sigma^2}= \frac{0,30197}{10} = 0,03020
			\]

			Entretanto, vamos preferir utilizar \(S^{2}\) como estimador da variância do erro, pois ele é um estimador não viesado.

		\end{solution}

		\vspace{0.3cm}

		\part

		\begin{solution}

			Para calcularmos a variância e covariância das variáveis em questão podemos utilizar uma abordagem matricial. A matriz de variância-covariância dos estimadores é dada por:
			\[VarCov = S^2 \times (X'X)^{-1}\]
			Como já obtivemos \(S^2\) na questão anterior, precisamos apenas calcular \((X'X)^{-1}\).

			\[(X'X) = \begin{bmatrix}
					n      & \sum X     \\
					\sum X & \sum X^{2}
				\end{bmatrix}
			\]

			E, portanto:
			\[(X'X)^{-1} = \frac{1}{n\sum X^{2} - (\sum X)^{2}} \times \begin{bmatrix}
					\sum X^{2} & -\sum X \\
					-\sum X    & n
				\end{bmatrix}
			\]

			Logo, temos que:
			\[
				VarCov = S^2 \times (X'X)^{-1} = 0,03775 \times \frac{1}{10 \times 7709 - (273)^{2}} \times
				\begin{bmatrix}
					7709 & -273 \\
					-273 & 10
				\end{bmatrix}
			\]
			\[
				VarCov =
				\begin{bmatrix}
					0,11363  & -0,00402 \\
					-0,00402 & 0,00015
				\end{bmatrix}
			\]

			Por fim, temos que:
			\[Var(\hat{\beta_{0}}) = 0,11363\]
			\[Var(\hat{\beta_{1}}) = 0,00015\]

		\end{solution}

		\vspace{0.3cm}

		\part

		\begin{solution}

			A hipótese que garante o não-viés dos estimadores é a hipótese de exogeneidade (H1), que diz que \(E(u|X) = 0\).

		\end{solution}

		\vspace{0.3cm}

		\part

		\begin{solution}

			A variância de \(Var(\hat{\beta_{0}} + \hat{\beta_{1}})\) é dada por:

			\[Var(\hat{\beta_{0}} + \hat{\beta_{1}}) = Var(\hat{\beta_{0}}) + Var(\hat{\beta_{1}}) + 2 \times Cov(\hat{\beta_{0}} \hat{\beta_{1}})\]
			Da letra C temos que:
			\[Var(\hat{\beta_{0}} + \hat{\beta_{1}}) = 0,11363 + 0,00015 + 2 \times (-0,00402)\]
			\[Var(\hat{\beta_{0}} + \hat{\beta_{1}}) = 0,10574\]

		\end{solution}

		\vspace{0.3cm}

		\part

		\begin{solution}

			Sabemos que \(R^{2}\) é dado por:
			\[ R^{2} = 1 - \frac{SQR}{SQT} \]

			E da letra b temos que \(SQR = 0,30197\). Agora precisamos calcular \(SQT\).

			\[SQT = \sum (Y-\bar{Y})^{2}\]
			\[ (3,0 - 3,41)^{2} =  0,1681 \]
			\[ (3,5 - 3,41)^{2} =  0,0081 \]
			\[ (3,6 - 3,41)^{2} =  0,0361 \]
			\[ (3,1 - 3,41)^{2} =  0,0961 \]
			\[ (3,3 - 3,41)^{2} =  0,0121 \]
			\[ (3,8 - 3,41)^{2} =  0,1521 \]
			\[ (4,0 - 3,41)^{2} =  0,3481 \]
			\[ (3,1 - 3,41)^{2} =  0,0961 \]
			\[ (3,2 - 3,41)^{2} =  0,0441 \]
			\[ (3,5 - 3,41)^{2} =  0,0081 \]
			\[SQT = 0,9690\]

			Portanto: \text{\large\( R^{2} = 1 - \frac{0,30197}{0,9690} = 0,68837\)}

		\end{solution}

		\vspace{0.3cm}

		\part

		\begin{solution}

			O modelo calculado foi:
			\[ retorno_{i} = 2,01688 + 0,05103 \times BTD_{i} + u_{i} \]

			Substituindo o valor apropriado de BTD, temos:
			\[ retorno_{i} = 2,01688 + 0,05103 \times 29 = 3,49675 \]

		\end{solution}
	\end{parts}

	\question Conforme Exercício 2

	\begin{parts}

		\part

		\begin{solution}

			\begin{center}
				\begin{tabular}{|c|c|}
					\hline
					Preço (Y) & Receita Líquida (X) \\ \hline
					10        & 25                  \\ \hline
					15        & 30                  \\ \hline
					30        & 35                  \\ \hline
					16        & 28                  \\ \hline
					20        & 25                  \\ \hline
					12        & 25                  \\ \hline
				\end{tabular}
			\end{center}

			Utilizaremos uma abordagem matricial para derivar os estimadores de MQO. As matrizes \(X\) e \(Y\) são dadas por:
			\[ X =
				\begin{bmatrix}
					1 & 25 \\
					1 & 30 \\
					1 & 35 \\
					1 & 28 \\
					1 & 25 \\
					1 & 25
				\end{bmatrix} \quad
				Y =
				\begin{bmatrix}
					10 \\
					15 \\
					30 \\
					16 \\
					20 \\
					12
				\end{bmatrix}
			\]

			Também podemos escrever a matriz \(X'\) da seguinte forma:
			\[
				X' =
				\begin{bmatrix}
					1  & 1  & 1  & 1  & 1  & 1  \\
					25 & 30 & 35 & 28 & 25 & 25
				\end{bmatrix}
			\]

			Logo, multiplicando \(X'\) por \(X\), obtemos:
			\[
				X'X =
				\begin{bmatrix}
					6   & 168  \\
					168 & 4784
				\end{bmatrix}
			\]

			Invertendo a matriz \(X'X\) obtemos:
			\[
				(X'X)^{-1} =
				\frac{1}{6 \times 4784 - 168^2} \times \begin{bmatrix}
					4784 & -168 \\
					-168 & 6
				\end{bmatrix}
			\]

			Podemos também obter a matriz \(X'Y\), que é dada por:
			\[
				X'Y =
				\begin{bmatrix}
					103 \\
					2998
				\end{bmatrix}
			\]

			Por fim, obtemos o produto de \((X'X)^{-1}\) por \(X'Y\):
			\[
				(X'X)^{-1}X'Y = \hat{\beta} =
				\frac{1}{6 \times 4784 - 168^2} \times \begin{bmatrix}
					-10912 \\
					684
				\end{bmatrix} =
				\begin{bmatrix}
					-22,73333 \\
					1,425
				\end{bmatrix} =
				\begin{bmatrix}
					\hat{\beta_{0}} \\
					\hat{\beta_{1}}
				\end{bmatrix}
			\]

			Logo, podemos interpretar que, quando a receita líquida é igual a 0, o preço do ativo é, em média, de R\$ -22,73333. Além disso, quando a receita líquida aumenta em (1) mil reais, o preço do ativo aumenta, em média, em R\$ 1,425.

		\end{solution}

		\vspace{0.3cm}

		\part

		\begin{solution}

			A variância dos estimadores pode ser calculada utilizando a seguinte fórmula:
			\[VarCov = S^2 \times (X'X)^{-1}\]

			Primeiramente precisamos calcular \(S^2\). Para isso, precisamos calcular o somatório do quadrado dos resíduos (SQR).
			\[SQR = \sum (Y - \widehat{Y})^{2}\]
			\[ [ 10 - (\hat{\beta_{0}} + \hat{\beta_{1}} \times 25)]^{2} =  8,36176 \]
			\[ [ 15 - (\hat{\beta_{0}} + \hat{\beta_{1}} \times 30)]^{2} =  25,16698 \]
			\[ [ 30 - (\hat{\beta_{0}} + \hat{\beta_{1}} \times 35)]^{2} =  8,17005 \]
			\[ [ 16 - (\hat{\beta_{0}} + \hat{\beta_{1}} \times 28)]^{2} =  1,36112 \]
			\[ [ 20 - (\hat{\beta_{0}} + \hat{\beta_{1}} \times 25)]^{2} =  50,52836 \]
			\[ [ 12 - (\hat{\beta_{0}} + \hat{\beta_{1}} \times 25)]^{2} =  0,79508 \]
			\[ SQR = 94,38335\]

			Logo, sabendo que \text{\large\(S^2 = \frac{SQR}{n-k}\)}:
			\[S^{2} = \frac{94,38335}{6-2} = 23,59584\]

			Agora podemos calcular a matriz \(VarCov\):
			\[VarCov = 23,59584 \times (X'X)^{-1}\]
			\[
				VarCov =  \frac{23,59584}{6 \times 4784 - 168^2} \times
				\begin{bmatrix}
					4784 & -168 \\
					-168 & 6
				\end{bmatrix}
			\]
			\[
				VarCov =
				\begin{bmatrix}
					235,17187 & -8,25854 \\
					-8,25854  & 0,29495
				\end{bmatrix}
			\]

			Portanto, temos:
			\[Var(\hat{\beta_{0}}) = 235,17187\]
			\[Var(\hat{\beta_{1}}) = 0,29495\]

		\end{solution}

		\vspace{0.3cm}

		\part

		\begin{solution}

			\begin{center}
				\begin{tabular}{|c|c|c|}
					\hline
					Preço (Y) & Receita Líquida (\(X_{1}\)) & Escolaridade (\(X_{2}\)) \\ \hline
					10        & 25                          & 22                       \\ \hline
					15        & 30                          & 20                       \\ \hline
					30        & 35                          & 24                       \\ \hline
					16        & 28                          & 16                       \\ \hline
					20        & 25                          & 18                       \\ \hline
					12        & 25                          & 24                       \\ \hline
				\end{tabular}
			\end{center}

			Por se tratar de uma regressão múltipla, precisaremos novamente adotar uma abordagem matricial. As matrizes relevantes são:
			\[
				X =
				\begin{bmatrix}
					1 & 25 & 22 \\
					1 & 30 & 20 \\
					1 & 35 & 24 \\
					1 & 28 & 16 \\
					1 & 25 & 18 \\
					1 & 25 & 24
				\end{bmatrix} \quad
				Y =
				\begin{bmatrix}
					10 \\
					15 \\
					30 \\
					16 \\
					20 \\
					12
				\end{bmatrix} \quad
				X' =
				\begin{bmatrix}
					1  & 1  & 1  & 1  & 1  & 1  \\
					25 & 30 & 35 & 28 & 25 & 25 \\
					22 & 20 & 24 & 16 & 18 & 24
				\end{bmatrix}
			\]

			Primeiramente obtemos \(X'X\):
			\[
				X'X =
				\begin{bmatrix}
					6   & 168  & 124  \\
					168 & 4784 & 3488 \\
					124 & 3488 & 2616
				\end{bmatrix}
			\]

			Em seguida obtemos também o determinante da matriz \(X'X\), que é igual a:
			\[det(X'X) = 24064\]

			Precisamos agora obter a matriz adjunta, que é simplesmente a transposta da matriz dos cofatores. Vamos agora calcular cada um dos cofatores.

			\[ C_{11} = -1^{2} \times det
				\begin{bmatrix}
					4784 & 3488 \\
					3488 & 2616
				\end{bmatrix}
				= 348800
				\qquad
				C_{12} = -1^{3} \times det
				\begin{bmatrix}
					168 & 3488 \\
					124 & 2616
				\end{bmatrix}
				= -6976
			\]

			\[ C_{13} = -1^{4} \times det
				\begin{bmatrix}
					168 & 4784 \\
					124 & 3488
				\end{bmatrix}
				= -7232
				\qquad
				C_{22} = -1^{4} \times det
				\begin{bmatrix}
					6   & 124  \\
					124 & 2616
				\end{bmatrix}
				= 320
			\]

			\[ C_{23} = -1^{5} \times det
				\begin{bmatrix}
					6   & 168  \\
					124 & 3488
				\end{bmatrix}
				= -96
				\qquad
				C_{33} = -1^{6} \times det
				\begin{bmatrix}
					6   & 168  \\
					168 & 4784
				\end{bmatrix}
				= 480
			\]

			Assim obtemos a matriz  dos cofatores:

			\[ Cof(X'X) =
				\begin{bmatrix}
					348800 & -6976 & -7232 \\
					-6976  & 320   & -96   \\
					-7232  & -96   & 480
				\end{bmatrix}
			\]

			Dessa forma podemos obter a matriz inversa de \(X'X\) da seguinte forma:
			\[(X'X)^{-1} = \frac{1}{det(X'X)} \times Adj(X'X)
			\]

			Onde \(Adj(X'X)\) é a matriz transposta da matriz dos cofatores. Como a matriz dos cofatores é simétrica, a matriz adjunta é igual a matriz dos cofatores.

			Ou seja:
			\[(X'X)^{-1} =
				\frac{1}{24064} \times
				\begin{bmatrix}
					348800 & -6976 & -7232 \\
					-6976  & 320   & -96   \\
					-7232  & -96   & 480
				\end{bmatrix}
			\]

			Obtemos agora a matriz \(X'Y\):
			\[
				X'Y =
				\begin{bmatrix}
					103  \\
					2998 \\
					2144
				\end{bmatrix}
			\]

			Podemos então calcular os coeficientes \(\hat{\beta}\):
			\[(X'X)^{-1}X'Y = \hat{\beta} =
				\frac{1}{24064} \times
				\begin{bmatrix}
					348800 & -6976 & -7232 \\
					-6976  & 320   & -96   \\
					-7232  & -96   & 480
				\end{bmatrix} \times
				\begin{bmatrix}
					103  \\
					2998 \\
					2144
				\end{bmatrix}
			\]
			\[(X'X)^{-1}X'Y = \hat{\beta} =
				\frac{1}{24064} \times
				\begin{bmatrix}
					-493056 \\
					35008   \\
					-3584
				\end{bmatrix} =
				\begin{bmatrix}
					-20,48936 \\
					1,45479   \\
					-0,14894
				\end{bmatrix} =
				\begin{bmatrix}
					\hat{\beta_{0}} \\
					\hat{\beta_{1}} \\
					\hat{\beta_{2}}
				\end{bmatrix}
			\]

		\end{solution}

		\vspace{0.3cm}

		\part

		\begin{solution}

			A literatura sugere uma relação positiva entre a escolaridade e o preço do ativo, ou seja, espera-se que \(\hat{\beta_{2}} > 0\). Entretanto, o valor obtido para \(\hat{\beta_{2}}\) foi \(-0,14894\), indicando uma relação negativa entre as variáveis. Existem evidências de que a hipótese de exogeneidade foi violada, uma vez que \(\hat{\beta}^{simples}_{1} \neq \hat{\beta}^{multipla}_{1}\).

		\end{solution}

		\vspace{0.3cm}

		\part

		\begin{solution}

			Para calcularmos a variância dos estimadores, precisamos novamente calcular \(S^2\). Para isso, precisamos calcular o somatório do quadrado dos resíduos (SQR).
			\[SQR = \sum (Y - \widehat{Y})^{2}\]
			\[ [ 10 - (\hat{\beta_{0}} + \hat{\beta_{1}} \times 25 + \hat{\beta_{2}} \times 22)]^{2} =  6,77930 \]
			\[ [ 15 - (\hat{\beta_{0}} + \hat{\beta_{1}} \times 30 + \hat{\beta_{2}} \times 20)]^{2} =  26,78621 \]
			\[ [ 30 - (\hat{\beta_{0}} + \hat{\beta_{1}} \times 35 + \hat{\beta_{2}} \times 24)]^{2} =  9,89901 \]
			\[ [ 16 - (\hat{\beta_{0}} + \hat{\beta_{1}} \times 28 + \hat{\beta_{2}} \times 16)]^{2} =  3,46600 \]
			\[ [ 20 - (\hat{\beta_{0}} + \hat{\beta_{1}} \times 25 + \hat{\beta_{2}} \times 18)]^{2} =  46,24721 \]
			\[ [ 12 - (\hat{\beta_{0}} + \hat{\beta_{1}} \times 25 + \hat{\beta_{2}} \times 24)]^{2} =  0,09353 \]
			\[ SQR = 93,27126\]

			Logo, sabendo que \text{\large\(S^2 = \frac{SQR}{n-k}\)}:
			\[S^{2} = \frac{93,27126}{6-3} = 31,09042\]

			Sabendo também que \(VarCov = S^2 \times (X'X)^{-1}\), temos que:
			\[VarCov = 31,09042 \times (X'X)^{-1}\]
			\[VarCov =
				\frac{31,09042}{24064} \times
				\begin{bmatrix}
					348800 & -6976 & -7232 \\
					-6976  & 320   & -96   \\
					-7232  & -96   & 480
				\end{bmatrix}
			\]
			\[
				VarCov =
				\begin{bmatrix}
					450,64572 & -9,01291 & -9,34366 \\
					-9,01291  & 0,41344  & -0,12403 \\
					-9,34366  & -0,12403 & 0,62015
				\end{bmatrix}
			\]

			Portanto, temos:
			\[Var(\hat{\beta_{0}}) = 450,64572\]
			\[Var(\hat{\beta_{1}}) = 0,41344\]
			\[Var(\hat{\beta_{2}}) = 0,62015\]

		\end{solution}

	\end{parts}

	\question Conforme Exercício 2

	\begin{solution}
		Inicialmente utilizamos nossos conhecimentos de regressão simples e sabemos que \text{\large\(\hat{\beta_{1}} = \frac{n \sum XY - \sum X \sum Y}{n \sum X^{2} - (\sum X)^{2}}\)} e que \(\hat{\beta_{0}} = \bar{Y} - \hat{\beta_{1}} \times \bar{X}\).

		Agora precisamos mostrar como calcular os os estimadores de MQO utilizando uma abordagem matricial. Sabendo que \(\sum{u_{i}^{2}} = u'u\), podemos escrever o seguinte:

		\[min\sum{u_{1}^{2}} = min\: u'u\]

		Como \(u' = (Y - X\beta)'\) e \(u = (Y - X\beta)\) temos:

		\[S(\beta) = (Y - X\beta)'(Y - X\beta)\]

		Aplicando a propriedade distributiva:

		\[S(\beta) = Y'Y - Y'X\beta - \beta'X'Y + \beta'X'X\beta\]

		Note que: \(Y'X\beta = \beta'X'Y\). Logo:

		\[S(\beta) = Y'Y - 2\beta'X'Y + \beta'X'X\beta\]

		Resolvendo a condição de primeira ordem, ou seja, para encontrarmos o mínimo local desta função igualamos sua primeira derivada parcial à zero:

		\[\frac{\partial S(\beta)}{\partial \beta'} = 0\]
		\[\frac{\partial S(\beta)}{\partial \beta'} = -2X'Y + 2X'X\beta = 0\]
		\[-2X'Y + 2X'X\beta = 0\]
		\[-2X'Y = -2X'X\beta\]
		\[X'Y = X'X\beta\]
		\[X'X\hat{\beta} = X'Y\]

		Podemos multiplicar os dois lados desta equação por \((X'X)^{-1}\) para obtermos o seguinte:

		\[(X'X)^{-1}(X'X)\hat{\beta} = (X'X)^{-1}X'Y\]

		Como \((X'X)^{-1}(X'X) = I\), temos:

		\[I\hat{\beta} = (X'X)^{-1}X'Y\]
		\[\hat{\beta} = (X'X)^{-1}X'Y\]

		Desta abordagem matricial temos as seguintes matrizes para calcular os estimadores de interesse:
		\[ X =
			\begin{bmatrix}
				1   & X_{1} \\
				1   & X_{2} \\
				... & ...   \\
				1   & X_{n}
			\end{bmatrix} \quad
			Y =
			\begin{bmatrix}
				Y_{1} \\
				Y_{2} \\
				...   \\
				Y_{n}
			\end{bmatrix} \quad
			X' =
			\begin{bmatrix}
				1     & 1     & ... & 1     \\
				X_{1} & X_{2} & ... & X_{n}
			\end{bmatrix}
		\]

		Podemos obter, então, a matriz \(X'X\):
		\[  X'X =
			\begin{bmatrix}
				n      & \sum X     \\
				\sum X & \sum X^{2}
			\end{bmatrix}
		\]
		Invertendo a matriz \(X'X\) obtemos:
		\[(X'X)^{-1} = \frac{1}{n \sum X^{2} - (\sum X)^{2}} \times
			\begin{bmatrix}
				\sum X^{2} & -\sum X \\
				-\sum X    & n
			\end{bmatrix}
		\]
		Obtemos também a matriz \(X'Y\):
		\[  X'Y =
			\begin{bmatrix}
				\sum Y \\
				\sum XY
			\end{bmatrix}
		\]

		Por fim, obtemos o produto de \((X'X)^{-1}\) por \(X'Y\):
		\[(X'X)^{-1}X'Y = \hat{\beta} =
			\frac{1}{n \sum X^{2} - (\sum X)^{2}} \times
			\begin{bmatrix}
				\sum X^{2}\sum Y - \sum X\sum XY \\
				n\sum XY - \sum X\sum Y
			\end{bmatrix} =
			\begin{bmatrix}
				\hat{\beta_{0}} \\
				\hat{\beta_{1}}
			\end{bmatrix}
		\]

		\[\hat{\beta_{0}} = \frac{\sum X^{2}\sum Y - \sum X\sum XY}{n \sum X^{2} - (\sum X)^{2}}\]
		\[\hat{\beta_{1}} = \frac{n\sum XY - \sum X\sum Y}{n \sum X^{2} - (\sum X)^{2}}\]

		O coeficiente \(\hat{\beta_{1}}\) calculado pela abordagem matricial concorda com o valor conhecido de \(\hat{\beta_{1}}\). Entretanto precisamos mostrar que o valor obtido de \(\hat{\beta_{0}}\) também é equivalente ao valor conhecido. Podemos adotar uma prova por contradição, para fazer isso escreveremos que o \(\hat{\beta_{0}}\) obtido pela via matricial é igual ao valor conhecido de \(\hat{\beta_{0}}\).

		\[ \frac{\sum X^{2}\sum Y - \sum X\sum XY}{n\sum X^{2} - (\sum X)^{2}} = \bar{Y} - \hat{\beta_{1}}\bar{X} \]

		Escrevendo \(\bar{Y}\) e \(\bar{X}\) em termos de somatórios e substituindo \(\hat{\beta_{1}}\) pela sua forma conhecida, temos:


		\[ \frac{\sum X^{2}\sum Y - \sum X\sum XY}{n\sum X^{2} - (\sum X)^{2}} = \frac{\sum Y}{n} - \frac{n\sum XY - \sum X\sum Y}{n\sum X^{2} - (\sum X)^{2}} \times  \frac{\sum X}{n}\]
		\[ \frac{\sum X^{2}\sum Y - \sum X\sum XY}{n\sum X^{2} - (\sum X)^{2}} + \frac{n \sum XY - \sum X\sum Y}{n\sum X^{2} - (\sum X)^{2}} \times \frac{\sum X}{n} = \frac{\sum Y}{n} \]
		\[ \frac{\sum X^{2}\sum Y - \sum X\sum XY}{n\sum X^{2} - (\sum X)^{2}} + \frac{\sum X(n\sum XY - \sum X\sum Y)}{n(n \sum X^{2} - (\sum X)^{2})} = \frac{\sum Y}{n} \]
		\[ \frac{n(\sum X^{2}\sum Y - \sum X\sum XY) + \sum X(n\sum XY - \sum X\sum Y)}{n(n \sum X^{2} - (\sum X)^{2})} = \frac{\sum Y}{n} \]
		\[ \frac{n(\sum X^{2}\sum Y - \sum X\sum XY) + n\sum X\sum XY - \sum X\sum X\sum Y}{n(n \sum X^{2} - (\sum X)^{2})} = \frac{\sum Y}{n} \]
		\[ \frac{n(\sum X^{2}\sum Y - \sum X\sum XY + \sum X\sum XY) - (\sum X)^{2}  \sum Y}{n(n \sum X^{2} - (\sum X)^{2})} = \frac{\sum Y}{n} \]
		\[ \frac{n\sum X^{2}\sum Y - (\sum X)^{2}  \sum Y}{n(n\sum X^{2} - (\sum X)^{2})} = \frac{\sum Y}{n} \]
		\[ \frac{n\sum X^{2}\sum Y - (\sum X)^{2}  \sum Y}{n\sum X^{2} - (\sum X)^{2}} = \sum Y \]
		\[ n\sum X^{2}\sum Y - (\sum X)^{2}\sum Y = \sum Y(n \sum X^{2} - (\sum X)^{2})\]
		\[ n\sum X^{2}\sum Y - (\sum X)^{2}\sum Y = n\sum Y\sum X^{2} - \sum Y(\sum X)^{2}\]
		\[ n\sum X^{2}\sum Y - (\sum X)^{2}\sum Y = n\sum Y\sum X^{2} - (\sum X)^{2}\sum Y\]
		\[ n\sum X^{2}\sum Y = n\sum Y\sum X^{2} \]
		\[ n\sum X^{2}\sum Y = n\sum X^{2}\sum Y \]
		\[ n = n\]

		Portanto, o valor de \(\hat{\beta_{0}}\) obtido pela via matricial é equivalente ao valor conhecido de \(\hat{\beta_{0}}\).

		Precisamos agora mostrar que, ao atendermos H1, os valores dos coeficiente obtidos são não viesados. Para esta condição ser satisfeita precisamos que \(\mathbb{E}(\hat{\beta}) = \beta\)

		Sabemos que:
		\[\hat{\beta} = (X'X)^{-1}X'Y\]

		E sabemos que:
		\[Y = X\beta + u\]

		Substituindo \(Y\) na equação de \(\hat{\beta}\), temos:
		\[\hat{\beta} = (X'X)^{-1}X'(X\beta + u)\]
		\[\hat{\beta} = (X'X)^{-1}X'X\beta + (X'X)^{-1}X'u\]
		\[\hat{\beta} = I\beta + (X'X)^{-1}X'u\]
		\[\hat{\beta} = \beta + (X'X)^{-1}X'u\]
		\[\mathbb{E}(\hat{\beta} \mid X) = \mathbb{E}(\beta + (X'X)^{-1}X'u \mid X)\]
		\[\mathbb{E}(\hat{\beta} \mid X) = \mathbb{E}(\beta \mid X) + \mathbb{E}((X'X)^{-1}X' \mid X) \mathbb{E}(u \mid X)\]

		Pela hipótese de exogeneidade sabemos que \(\mathbb{E}(u \mid X) = 0\), logo:
		\[\mathbb{E}(\hat{\beta} \mid X) = \mathbb{E}(\beta \mid X) + \mathbb{E}((X'X)^{-1}X' \mid X) \times 0\]
		\[\mathbb{E}(\hat{\beta} \mid X) = \mathbb{E}(\beta \mid X)\]
		\[\mathbb{E}(\hat{\beta} \mid X) = \beta\]

	\end{solution}

\end{questions}
\end{document}
